% !TeX root = qa_report.tex
% !TeX program = pdflatex
\documentclass[11pt,a4paper]{article}

\pdfminorversion=7

% =====================
% Packages (IDENTICAL TO QC)
% =====================
\usepackage[top=0.9in,left=0.9in,right=0.9in,bottom=1.2in]{geometry}
\usepackage{array}
\usepackage{ragged2e}
\usepackage{setspace}
\usepackage{hyperref}
\usepackage{enumitem}
\usepackage{graphicx}
\usepackage{fancyhdr}
\usepackage{tabularx}
\usepackage{booktabs}
\usepackage{multirow}
\usepackage{float}
\usepackage{xcolor}
\usepackage{parskip}
\usepackage{amsmath}
\usepackage{amssymb}
\usepackage{colortbl}
\usepackage{longtable}
\usepackage{textcomp}
\usepackage{caption}
\usepackage{bookmark}
\usepackage[normalem]{ulem}

\hypersetup{
colorlinks=true,
linkcolor=blue,
urlcolor=blue,
citecolor=blue
}

\AtBeginDocument{
\let\origref\ref
\renewcommand{\ref}[1]{\uline{\origref{#1}}}
\let\orighyperref\hyperref
\renewcommand{\hyperref}[2][]{\orighyperref[#1]{\uline{#2}}}
}

\definecolor{tableheader}{RGB}{240,240,240}
\definecolor{tablerow}{RGB}{250,250,250}

\frenchspacing

% =====================
% Font
% =====================
\usepackage[T1]{fontenc}
\usepackage[scaled=0.95]{helvet}
\renewcommand{\familydefault}{\sfdefault}

\setstretch{1.0}

% =====================
% BACKGROUND PLACEHOLDER
% =====================
% Replace later with the same background setup used in SRS/QC

% =====================
% Header/Footer
% =====================
\pagestyle{fancy}
\fancyhf{}
\lhead{\textbf{[QA AUDIT ORGANISATION NAME]}}
\rhead{Quality Assurance (QA) Audit Report}
\cfoot{\thepage}
\renewcommand{\headrulewidth}{0.4pt}
\setlength{\headheight}{20pt}

\begin{document}

% =====================
% COVER PAGE
% =====================
\thispagestyle{empty}
\fancyhf{}
\renewcommand{\headrulewidth}{0pt}

\newgeometry{top=5.2cm,left=0.9in,right=0.9in,bottom=1.2in}

\begin{center}

{\Large\fbox{\parbox{0.72\textwidth}{\centering
\vspace{0.3cm}
\textbf{BACKGROUND PLACEHOLDER}\\
Insert QA cover background here later
\vspace{0.3cm}}}}

\vspace{1cm}

\LARGE\textbf{Quality Assurance (QA) Audit Report}

\large Requirements Engineering Process

\large Abyss Navigator Suite Project

\vspace{0.8cm}

\textbf{\Large [QA AUDIT ORGANISATION NAME]}

\today

\end{center}

\begin{table}[H]
\centering
\small
\renewcommand{\arraystretch}{2.0}
\begin{tabular}{|m{0.25\textwidth}|m{0.2\textwidth}|m{0.25\textwidth}|m{0.2\textwidth}|}
\hline
\rowcolor{tableheader}
QA Team Member & Department & Role & Employee ID \\
\hline
Zong Han & QA Audit Division & QA Auditor & QA-001 \\
\hline
\end{tabular}
\caption{QA Audit Personnel}
\label{tab:qa-team}
\end{table}

\restoregeometry
\newpage

% =====================
% Table of Contents
% =====================
\renewcommand{\contentsname}{Table of Contents}
\thispagestyle{empty}

\hypersetup{linkcolor=black}
\tableofcontents
\newpage
\hypersetup{linkcolor=blue}

\noindent\fbox{\begin{minipage}{\dimexpr\textwidth-2\fboxsep-2\fboxrule\relax}
\small\textbf{Note on Hyperlinks:} Throughout this document, all hyperlinks are displayed in \textcolor{blue}{\underline{blue and underlined}}. These links are fully clickable in the PDF viewer and will navigate to the relevant section, finding, or appendix.
\end{minipage}}

\bigskip

% =====================
% EXECUTIVE SUMMARY
% =====================

\section{Executive Summary}

This report presents the results of a Quality Assurance (QA) audit of the
Requirements Engineering (RE(-)) process used to produce the Software
Requirements Specification (SRS) for the \textbf{Abyss Navigator Suite} project.

The purpose of this QA audit is to determine whether the vendor's conduct of
Analysing, Specifying, and Validating followed disciplined and sufficiently
documented Requirements Engineering practices. In accordance with the
Validating method, QA targets the \textbf{process that produces the SRS},
rather than the SRS work product itself.

The audit focused on process non-compliance risks such as incomplete process
evidence, weak traceability discipline, insufficient documentation of review
activities, and uneven coverage of method-specific artefacts. While the RE(-)
process shows a generally structured and credible approach, several areas
require strengthening to improve auditability, process transparency, and
confidence in method compliance.

Overall, the QA audit concludes that the RE(-) process is \textbf{conditionally
acceptable}, subject to rectification of the identified process non-compliance
issues before future downstream governance and project-control activities rely
on the process record as complete and robust.

% =====================
% 1 INTRODUCTION
% =====================

\section{Introduction}

\subsection{Purpose}

This Quality Assurance (QA) Audit Report evaluates the Requirements Engineering (RE(-)) process used in producing the Software Requirements Specification (SRS) for the \textbf{Abyss Navigator Suite} project.

While Quality Control (QC) focuses on detecting defects within the SRS work product, Quality Assurance focuses on evaluating whether the \textbf{process used to produce the requirements followed appropriate Requirements Engineering practices}.

The purpose of this QA audit is therefore to determine whether the RE(-) process applied during the project followed disciplined and systematic procedures.

\subsection{Document Conventions}

This document is an audit report rather than a process manual. It records detected process non-compliance issues, explains their implications, and recommends rectification actions.

Error detection instances use the format \texttt{QA-ERR-XXX}.

Audit findings use the format \texttt{QA-FND-XXX}.

Rectification recommendations use the format \texttt{QA-REC-XXX}.

\subsection{Intended Audience and Reading Suggestions}

This report is intended for Customer C-staff, project management, the requirements engineering team, and other stakeholders responsible for process governance in Requirements Engineering.

Readers should begin with the Executive Summary and Sections 1--3, then proceed to the QA Checklist, Error Detection, detailed Findings, and Recommended Process Rectifications.

\subsection{Scope of QA Audit}

The QA audit focuses on evaluating the process used during the following RE(-) activities:

\begin{itemize}[leftmargin=*]
\item requirements analysis
\item requirements specification
\item requirements validation
\item requirements traceability management
\item requirements change-mitigation setup
\item process evidence and governance records
\end{itemize}

% =====================
% 2 AUDIT BASIS
% =====================

\section{Audit Basis}

\subsection{Nature of QA}

Quality Assurance audits the \textbf{process used to produce requirements}, rather than the SRS document itself.

The QA auditor therefore evaluates whether the vendor team followed structured Requirements Engineering practices during Analysing, Specifying, and Validating.

\subsection{Basis of Evaluation}

The QA audit was conducted on the basis that the vendor should have followed process expectations associated with:

\begin{itemize}[leftmargin=*]
\item regulated-revision during Analysing
\item documented incorporation of development provisions during Specifying
\item planned and report-based auditing during Validating
\item maintenance of supporting plans, reports, checklists, meeting records, and artefacts throughout RE(-)
\end{itemize}

\subsection{Kinds of Errors Pursued}

The QA audit investigates whether the Requirements Engineering process suffered from failure to strictly comply with established RE(-) practices.

Potential causes examined include:

\begin{itemize}[leftmargin=*]
\item insufficient or invalid process evidence
\item insufficient resources allocated to requirements activities
\item inadequate management tools supporting requirements work
\item insufficient training or guidance in applying RE(-) techniques
\item incomplete linkage between process tasks and artefacts
\end{itemize}

\subsection{Severity Classification}

Findings are categorised according to severity:

\begin{itemize}[leftmargin=*]
\item \textbf{High} -- process non-compliance that materially weakens confidence in whether RE(-) methods were properly conducted
\item \textbf{Medium} -- process weakness that reduces auditability, process clarity, or evidence quality
\item \textbf{Low} -- minor documentary or structural weakness that does not materially undermine overall process conduct
\end{itemize}

\subsection{Records Reviewed}

The following project records were examined during the QA audit:

\begin{itemize}[leftmargin=*]
\item Project Specification (PS)
\item Software Requirements Specification (SRS)
\item requirements modelling artefacts
\item requirement traceability records
\item meeting documentation and review records
\item project planning and documentation materials
\item supporting notes relevant to Analysing, Specifying, and Validating conduct
\end{itemize}

% =====================
% 3 QA PLAN
% =====================

\section{QA Plan}

\subsection{Target Process Description}

The target of this QA audit is the vendor's \textbf{Requirements Engineering process} used to transform the Project Specification into a validated SRS through Analysing, Specifying, and Validating.

\subsection{Context and Goals}

The RE(-) process is expected to convert customer-originated requirements content into a vendor-ready and validated SRS that can be relied upon by downstream development activities. Since the quality of the final SRS depends partly on the integrity of the process that produced it, QA is conducted to assess whether that process was sufficiently disciplined and auditable.

The goals of this QA audit are:

\begin{itemize}[leftmargin=*]
\item to determine whether the RE(-) process followed systematic and sufficiently documented practices
\item to determine whether sufficient evidence exists to support key requirements development activities
\item to identify process weaknesses that may affect confidence in the quality and integrity of the resulting SRS
\end{itemize}

\subsection{Access and Audit Conditions}

The QA auditor reviewed documentation and artefacts produced during the requirements engineering stages. Access was limited to materials necessary for evaluating the RE(-) process and its supporting records.

\subsection{Work Breakdown}

\begin{table}[H]
\centering
\small
\renewcommand{\arraystretch}{1.4}
\begin{tabularx}{\textwidth}{|p{0.24\textwidth}|p{0.22\textwidth}|X|}
\hline
\rowcolor{tableheader}
\textbf{Audit Role} & \textbf{Assigned Auditor} & \textbf{Primary Responsibility} \\
\hline
QA Auditor & Zong Han & Review process records, evaluate method compliance, identify process non-compliance instances, and prepare the final QA report for Customer C-staff. \\
\hline
\end{tabularx}
\caption{QA audit work allocation}
\label{tab:qa-wbs}
\end{table}

\subsection{Indicative Timeline}

\begin{table}[H]
\centering
\small
\renewcommand{\arraystretch}{1.4}
\begin{tabularx}{\textwidth}{|p{0.18\textwidth}|p{0.42\textwidth}|X|}
\hline
\rowcolor{tableheader}
\textbf{Stage} & \textbf{Activity} & \textbf{Indicative Duration} \\
\hline
Stage 1 & Review process documentation and supporting artefacts for Analysing, Specifying, and Validating & 2 hours \\
\hline
Stage 2 & Map observed records against expected RE(-) method conduct & 1 hour \\
\hline
Stage 3 & Identify non-compliance instances and analyse severity & 1 hour \\
\hline
Stage 4 & Finalise QA report and sign-off submission to Customer C-staff & 1 hour \\
\hline
\end{tabularx}
\caption{Indicative QA audit timeline}
\label{tab:qa-timeline}
\end{table}

\subsection{Reporting}

The QA auditor documents findings and recommendations in this report. The auditor does not directly modify the audited work products or process records.

% =====================
% 4 QA CHECKLIST
% =====================

\section{QA Checklist}

The following checklist was used to inspect whether the RE(-) process showed sufficient compliance, evidence, and process discipline.

\begin{table}[H]
\centering
\small
\renewcommand{\arraystretch}{1.4}
\begin{tabularx}{\textwidth}{|p{0.1\textwidth}|X|p{0.2\textwidth}|}
\hline
\rowcolor{tableheader}
\textbf{Item} & \textbf{Checklist Question} & \textbf{Inspection Result} \\
\hline
Q1 & Is there sufficient documented evidence that Analysing was carried out as regulated-revision rather than ad hoc rewriting? & Partially satisfied \\
\hline
Q2 & Is there evidence that Specifying incorporated development provisions such as quality model, iRTM, and change mitigation? & Mostly satisfied \\
\hline
Q3 & Is there evidence that Validating was conducted as an audit-oriented activity rather than informal review only? & Partially satisfied \\
\hline
Q4 & Are process records such as plans, notes, and artefacts sufficiently linked to the three RE(-) methods? & Partially satisfied \\
\hline
Q5 & Are traceability records and supporting models sufficiently governed as process artefacts? & Mostly satisfied \\
\hline
Q6 & Are assumptions, constraints, and method decisions explicitly recorded where needed? & Partially satisfied \\
\hline
Q7 & Is the process record sufficiently clear for an independent auditor to understand what was done and how? & Partially satisfied \\
\hline
Q8 & Are supporting addenda and referenced process artefacts appropriately linked to the base report? & Mostly satisfied \\
\hline
\end{tabularx}
\caption{QA process checklist}
\label{tab:qa-checklist}
\end{table}

% =====================
% 5 QA ERROR DETECTION
% =====================

\section{QA Error Detection}

The QA audit identified instances of \textbf{process non-compliance} in the vendor's conduct of Requirements Engineering. Each instance is recorded with its location and characteristic.

\begin{table}[H]
\centering
\small
\renewcommand{\arraystretch}{1.4}
\begin{tabularx}{\textwidth}{|p{0.15\textwidth}|p{0.22\textwidth}|p{0.2\textwidth}|X|}
\hline
\rowcolor{tableheader}
\textbf{Error ID} & \textbf{Error Type} & \textbf{Location} & \textbf{Characteristic of Error} \\
\hline
QA-ERR-001 & Process non-compliance & Analysing process evidence & Process records do not always make the regulated-revision trail fully explicit from issue identification to corrected requirement outcome. \\
\hline
QA-ERR-002 & Process non-compliance & Validation / review records & Documentation of formal stakeholder validation or clarification activity could be stronger in the retained process record. \\
\hline
QA-ERR-003 & Process non-compliance & Traceability governance & The process record shows traceability artefacts, but the governance linkage between source decisions and downstream traceability maintenance can be clearer. \\
\hline
QA-ERR-004 & Process non-compliance & Method documentation & Some assumptions, constraints, and supporting method decisions are present in artefacts but not always captured as explicit process evidence. \\
\hline
\end{tabularx}
\caption{Detected QA error instances}
\label{tab:qa-error-detection}
\end{table}

% =====================
% 6 QA FINDINGS
% =====================

\section{QA Findings}

\subsection{Finding Summary}

\begin{table}[H]
\centering
\small
\renewcommand{\arraystretch}{1.4}
\begin{tabularx}{\textwidth}{|p{0.15\textwidth}|p{0.22\textwidth}|p{0.15\textwidth}|X|}
\hline
\rowcolor{tableheader}
\textbf{Finding ID} & \textbf{Type} & \textbf{Severity} & \textbf{Short Description} \\
\hline
QA-FND-001 & Process evidence weakness & Medium & The regulated-revision trail for some analysed content is not always explicit enough for independent reconstruction of the process path. \\
\hline
QA-FND-002 & Review / validation evidence weakness & Medium & Stronger retained evidence of stakeholder validation and clarification activity would improve process confidence. \\
\hline
QA-FND-003 & Traceability governance weakness & Medium & Traceability is present, but the supporting process evidence for how it was governed can be made more explicit. \\
\hline
QA-FND-004 & Documentation discipline weakness & Low & Some assumptions, constraints, and method decisions would benefit from more explicit process-level recording. \\
\hline
\end{tabularx}
\caption{Summary of QA findings}
\label{tab:qa-findings-summary}
\end{table}

\subsection{Detailed Findings}

\subsubsection*{QA-FND-001: Regulated-revision trail not always explicit}
\textbf{Related error detection:} QA-ERR-001 \\
\textbf{Type:} Process evidence weakness \\
\textbf{Severity:} Medium

\textbf{Context.} Analysing in RE(-) is expected to function as regulated-revision of the Project Specification into ready-content for the SRS.

\textbf{Result.} The available process records support that analysis occurred, but some parts of the process trail are not always explicit enough for an independent reviewer to reconstruct the sequence from identified issue to final revised requirement state.

\textbf{Analysis.} This is a process non-compliance concern because the method may have been conducted substantively, yet the retained evidence does not always show the full disciplined pathway with sufficient visibility.

\textbf{Message.} Strengthen the retained record of regulated-revision activity by documenting issue identification, clarification input, interpretation outcome, and categorisation result more explicitly.

\subsubsection*{QA-FND-002: Stakeholder validation evidence can be stronger}
\textbf{Related error detection:} QA-ERR-002 \\
\textbf{Type:} Review / validation evidence weakness \\
\textbf{Severity:} Medium

\textbf{Context.} Requirements clarification and validation rely on cooperation between customer and vendor, especially where uncertainty or incongruity exists.

\textbf{Result.} The process artefacts suggest that validation and clarification occurred, but the formal retained evidence of those activities could be stronger and more consistently visible.

\textbf{Analysis.} This reduces confidence in strict process compliance because an independent auditor should be able to verify that required clarification and confirmation activities were actually carried out.

\textbf{Message.} Retain clearer meeting notes, review decisions, or other formal evidence showing how clarification and validation outcomes were achieved.

\subsubsection*{QA-FND-003: Traceability governance process needs clearer linkage}
\textbf{Related error detection:} QA-ERR-003 \\
\textbf{Type:} Traceability governance weakness \\
\textbf{Severity:} Medium

\textbf{Context.} Traceability is expected to operate as a governance mechanism linking requirements, artefacts, and downstream development records.

\textbf{Result.} Traceability artefacts exist, but the process record can more clearly show how traceability decisions, mappings, and maintenance responsibilities were governed.

\textbf{Analysis.} This is not a claim that traceability is absent. Rather, it is a process non-compliance concern that the supporting governance evidence is not as explicit as it could be.

\textbf{Message.} Add clearer process records showing how traceability mappings were established, checked, and maintained.

\subsubsection*{QA-FND-004: Some method decisions are insufficiently explicit as process evidence}
\textbf{Related error detection:} QA-ERR-004 \\
\textbf{Type:} Documentation discipline weakness \\
\textbf{Severity:} Low

\textbf{Context.} RE(-) process discipline benefits when assumptions, constraints, and method decisions are explicitly recorded rather than only inferable from artefacts.

\textbf{Result.} Some of these items appear to be represented in project materials, but not always in a form that makes the process trail explicit.

\textbf{Analysis.} This weakens process transparency slightly, although it does not on its own materially invalidate the overall RE(-) process.

\textbf{Message.} Record key assumptions, constraints, and process decisions more explicitly in retained process notes and supporting artefacts.

% =====================
% 7 RECOMMENDED PROCESS RECTIFICATIONS
% =====================

\section{Recommended Process Rectifications}

\subsection{Priority Rectification Actions}

Based on the process non-compliance issues identified during the QA audit, the following rectification actions are recommended in priority order:

\begin{enumerate}[leftmargin=*]
\item strengthen the retained regulated-revision trail for analysed requirements content
\item improve retained evidence of clarification, review, and validation activities
\item make traceability governance decisions and maintenance responsibilities more explicit
\item document assumptions, constraints, and method decisions more systematically
\end{enumerate}

\subsection{Recommended Ownership of Rectifications}

\begin{table}[H]
\centering
\small
\renewcommand{\arraystretch}{1.4}
\begin{tabularx}{\textwidth}{|p{0.16\textwidth}|p{0.28\textwidth}|X|}
\hline
\rowcolor{tableheader}
\textbf{Recommendation ID} & \textbf{Owner} & \textbf{Action} \\
\hline
QA-REC-001 & Requirements Engineering Lead & Strengthen process evidence linking identified issues, clarification outcomes, and revised requirements. \\
\hline
QA-REC-002 & Project Leader / Requirements Governance & Ensure validation and clarification sessions leave a stronger retained audit trail. \\
\hline
QA-REC-003 & Traceability Custodian / Reviewer & Record clearer governance of traceability mappings, reviews, and maintenance actions. \\
\hline
QA-REC-004 & Process Documentation Owner & Standardise recording of assumptions, constraints, and method decisions across RE(-) artefacts. \\
\hline
\end{tabularx}
\caption{Recommended ownership of QA rectification actions}
\label{tab:qa-recommendations}
\end{table}

% =====================
% 8 CONCLUSION
% =====================

\section{Conclusion}

The QA audit concludes that the Requirements Engineering process used in the Abyss Navigator Suite project demonstrates a generally structured approach to requirement development and validation.

However, several improvements are required in documentation discipline, retained process evidence, and governance clarity in order to strengthen confidence that the process strictly complied with the expected conduct of Analysing, Specifying, and Validating.

The QA auditor does not directly modify the audited process records. Instead, the findings and recommended process rectifications in this report should be used to improve the discipline, transparency, and auditability of the RE(-) process.

Overall, the RE(-) process is assessed as \textbf{conditionally acceptable subject to rectification of the identified process non-compliance issues}.

% =====================
% 9 SIGN-OFF
% =====================

\section{Sign-off}

\subsection{QA Audit Authentication}

\begin{table}[H]
\centering
\small
\renewcommand{\arraystretch}{1.6}
\begin{tabular}{|m{0.35\textwidth}|m{0.25\textwidth}|m{0.3\textwidth}|}
\hline
\rowcolor{tableheader}
\textbf{Name} & \textbf{Role} & \textbf{Authentication} \\
\hline
Zong Han & QA Auditor & Submitted by named auditor to Customer C-staff \\
\hline
\end{tabular}
\caption{QA audit sign-off}
\label{tab:qa-signoff}
\end{table}

% =====================
% 10 ADDENDA
% =====================

\section{Addenda}

\subsection{Repository Information}

\begin{center}
\fbox{\begin{minipage}{0.9\textwidth}
\centering
\textbf{PROJECT REPOSITORY}\\[2pt]
\href{https://github.com/Spooky312/The-Abyss-Navigation-Suite-INF2113}{\uline{\nolinkurl{https://github.com/Spooky312/The-Abyss-Navigation-Suite-INF2113}}}\\[4pt]
\textit{QA report, SRS, models, and supporting process artefacts maintained under version control}
\end{minipage}}
\end{center}

\subsection{Glossary}

\begin{table}[H]
\centering
\small
\renewcommand{\arraystretch}{1.4}
\begin{tabularx}{\textwidth}{|p{0.28\textwidth}|X|}
\hline
\rowcolor{tableheader}
\textbf{Term} & \textbf{Definition} \\
\hline
QA & Quality Assurance audit of the process used to produce the requirements work product. \\
\hline
RE(-) process & The vendor-side conduct of Requirements Engineering methods relevant to producing and validating the SRS. \\
\hline
Process non-compliance & A process weakness or deficiency indicating incomplete adherence to expected method conduct or supporting evidence. \\
\hline
Regulated-revision & Controlled revision of the Project Specification into improved requirements content for the SRS. \\
\hline
Traceability governance & The process discipline used to establish, maintain, and justify traceability links across requirements artefacts. \\
\hline
\end{tabularx}
\caption{QA glossary addendum}
\label{tab:qa-glossary}
\end{table}

% =====================
% References
% =====================
\clearpage
\pagenumbering{arabic}
\setcounter{page}{1}
\renewcommand{\thepage}{R-\arabic{page}}
\section*{References}
\addcontentsline{toc}{section}{References}
\begin{itemize}[leftmargin=*,nosep]
\item Triton Deepwater, \textit{Project Specification: The Abyss Navigator Suite}, ICT2113, 2026.
\item Brigman Deep Sea Technology, \textit{Software Requirements Specification (SRS): Abyss Navigator Suite}, 2026.
\item Kevin Anthony Jones, \textit{Labtorial Guide for Class Project: Reqrs now \& near future}, Dec 2025.
\item Kevin Anthony Jones, \textit{Rubric for Assessment of Validating Report}, Mar 2026.
\item Karl E. Wiegers, \textit{Software Requirements Specification Template}, 1999.
\end{itemize}

\end{document}